% sections/03_physics.tex: why the hand is detectable (mirror symmetry, chirality, the first-half yaw about gravity).
% The mirror derivation (eq:mirror, eq:parity) and tab:mirror come from the skeleton; the evidence follows.
\section{Why the hand is detectable}\label{sec:physics}

\subsection{The mirror}\label{sec:physics:mirror}
Let $\Rmat(t)$ map device coordinates to body coordinates, with $\vf$ the specific force (the
browser's acceleration including gravity), $\va$ the acceleration and $\vomega$ the angular
velocity, all in device coordinates. Mirroring the body across its sagittal plane, $\mathbf{p}' = \mathbf{M}\mathbf{p}$,
and relabelling the device's $x$ axis, $\Rmat' = \mathbf{M}\Rmat\Smat$ with
$\mathbf{M} = \Smat = \operatorname{diag}(-1, 1, 1)$, gives for every grip and every movement
\begin{equation}
  \vf' = \Smat\vf, \qquad \va' = \Smat\va, \qquad \vomega' = -\Smat\vomega .
  \label{eq:mirror}
\end{equation}
Acceleration is a polar vector; angular velocity is a pseudovector and picks up $\det\Smat = -1$.
With $\ghat$ the unit vector pointing up, the yaw rate about gravity, $\vomega\cdot\ghat$, is a pseudoscalar: a
scalar that changes sign under the mirror and is unchanged by any rotation. This report calls such quantities
chiral (mirror-odd), and quantities that keep their sign mirror-even:
\begin{equation}
  (-\Smat\vomega)\cdot(\Smat\ghat) = -\,\vomega\cdot\ghat
  \quad\text{(mirror)}, \qquad
  (\Qmat\vomega)\cdot(\Qmat\ghat) = \vomega\cdot\ghat
  \quad\text{(any grip } \Qmat \in SO(3)\text{)} .
  \label{eq:parity}
\end{equation}
Holding the phone with the screen facing away is the proper rotation
$\Qmat = \mathbf{R}_y(180^\circ) = \operatorname{diag}(-1, 1, -1)$, so it leaves every chiral feature
unchanged (Table~\ref{tab:mirror}). An antisymmetric model, with logit
$z(\mathbf{x}) = g(\mathbf{x}) - g(m(\mathbf{x}))$, where $m(\cdot)$ is the mirror, then satisfies
$\pright(m(\mathbf{x})) = 1 - \pright(\mathbf{x})$ exactly. Only chiral features can carry the hand: a mirror-even
feature takes the same distribution for both hands whenever the two hands move as mirror images.

\begin{table}[htbp]
\centering
\caption{The mirror and the screen-away flip in device coordinates.}\label{tab:mirror}
\footnotesize
\begin{tabular}{@{}lll@{}}
\toprule
Quantity & Mirror $m$ (left $\leftrightarrow$ right) & Screen away, $\Qmat = \mathbf{R}_y(180^\circ)$ \\
\midrule
$\vf$, $\va$, $\ghat$ (polar vectors) & $(-x, y, z)$ & $(-x, y, -z)$ \\
$\vomega$ (pseudovector) & $(x, -y, -z)$ & $(-x, y, -z)$ \\
$\vomega\cdot\ghat$, $\vomega\cdot\va$, $\det[\va(t), \va(t+\tau), \ghat]$ & sign flips & unchanged \\
$\lVert\vomega\rVert$, $\va\cdot\ghat$, duration (mirror-even) & unchanged & unchanged \\
Hand label & swapped & unchanged \\
Grip label & top left $\leftrightarrow$ top right & toward $\leftrightarrow$ away \\
\bottomrule
\end{tabular}
\end{table}

\subsection{The mechanism}\label{sec:physics:mechanism}
Early in the approach the forearm swings toward the body's midline. Seen from above, the phone therefore turns
counter-clockwise in the right hand and clockwise in the left hand. Near the nose the wrist turns the phone partly
back to face the person, so the net rotation about gravity over the whole approach is a weaker and more variable
signal than the early turn. The early turn carries the clearest signal.

\subsection{The evidence on 20 persons}\label{sec:physics:evidence}
The yaw about gravity accumulated over the first half of the approach separates the hands with a pooled area under
the receiver operating characteristic curve (AUC) of \R{describe/fhy_auc_pooled}. The net yaw over the whole
approach has a lower pooled AUC, \R{describe/net_yaw_auc_pooled}; the difference was not tested in a paired
comparison (Figure~\ref{fig:S2_13_first_half_yaw}). Within persons the median AUC is
\R{describe/fhy_person_auc_median}, and it is above 0.5 for \R{describe/fhy_persons_auc_above_half} of
\R{describe/n_persons} persons. The sign of the first-half yaw alone, with no training, classifies
\R{describe/fhy_sign_rule_accuracy} of the trials, and it separates the hands with the same sign on both platforms
(AUC \R{describe/fhy_auc_android} on Android, \R{describe/fhy_auc_ios} on iOS).

\begin{figure}[htbp]
\centering
\genfig{core/S2_13_first_half_yaw}
\caption{The yaw about gravity over the first half of the approach has a higher pooled AUC than the net yaw over
the whole approach (\R{describe/fhy_auc_pooled} against \R{describe/net_yaw_auc_pooled}; no paired test). Left:
AUC within each of the \R{describe/n_persons} persons; right: the pooled distributions. The early turn carries the
clearest signal, so the models get time-resolved inputs.}
\label{fig:S2_13_first_half_yaw}
\end{figure}

The class means tell the same story (Figure~\ref{fig:S2_11_class_means}). The yaw rate about gravity has opposite
signs for the two hands over most of the first half of the approach, and the cumulative yaw keeps most of that
early turn. The mirrored left-hand mean has the shape of the right-hand mean but a smaller early peak, and the
mirror-even channels differ slightly between the hands: the two hands are close to, but not exactly, mirror images.
Mirror augmentation is therefore a plausible, not an exact, way to share data between the hands. In the feature catalogue of Appendix~\ref{app:m1:catalogue}, chiral features
reach a LOSO AUC of up to \R{m1/cat_odd_auc_max}, while every mirror-even feature stays between
\R{m1/cat_even_auc_min} and \R{m1/cat_even_auc_max}. The sensor conventions behind the mirror are verified on the
data: of the 64 possible sign patterns, the mirror predicted by the conventions maps right-hand onto left-hand
waveforms best, on both platforms (Appendix~\ref{app:data:conventions}).

\begin{figure}[htbp]
\centering
\genfig{core/S2_11_class_means}
\caption{Mean time course of the gravity-frame channels over the time-normalised approach, left against right
hand (band: $\pm$ standard error across the \R{describe/n_persons} persons). Dashed: the left-hand mean mirrored,
which should lie on the right-hand mean if the movement is mirror-symmetric. The chiral channels carry the hand
early in the approach; the mirror-even channels do not.}
\label{fig:S2_11_class_means}
\end{figure}

\subsection{Consequences for modelling}\label{sec:physics:consequences}
Three consequences shape the models. First, they get chiral, time-resolved inputs in the gravity frame, not only
whole-movement integrals. Second, a body frame that follows the phone's heading gives chiral channels that do not
depend on the grip (Section~\ref{sec:methods}), which matters when the test grip was never seen in training
(Section~\ref{sec:results:robustness}). Third, the mirror doubles the training data: a mirrored trial is a valid
trial of the other hand (mirror augmentation), or the model can be built antisymmetric. The return of the arm is a
time-reversed approach, so its chiral signs are reversed: on the \R{describe/return_trials} trials with a clean
return window, the first-half yaw rule is right on \R{describe/return_rule_acc_approach} of the approaches and on
only \R{describe/return_rule_acc_return} of the returns.
A cue-free detector that takes the return for the approach therefore answers the wrong hand
(Appendix~\ref{app:prod:return}).

\subsection{Limits}\label{sec:physics:limits}
The mirror of the accelerometer is approximate, because the sensor sits off the phone's centre line. The two arms
need not move as exact mirror images either: the dominant and the non-dominant arm differ in how they are
controlled \cite{sainburg2002dominance,przybyla2012dominance}. Section~\ref{sec:secondary} finds no detectable
dominant-hand signal in mirror-invariant kinematics.
